\section{答题时间整体规划与答题技巧}

\subsection{试卷结构速览}

\begin{center}
\begin{tabular}{c|c|c|c}
\textbf{题型} & \textbf{题量} & \textbf{分值} & \textbf{建议用时} \\ \hline
选择题（必考） & 14 题 & 42 分 & 25 分钟 \\
非选择题（必考） & 3 题 & 37 分 & 35 分钟 \\
非选择题（选考，二选一） & 1 题 & 15 分 & 15 分钟 \\ \hline
总计 & 18 题 & 100 分 & 75 分钟
\end{tabular}
\end{center}

必考非选择题三大题型：
\begin{itemize}
  \item 工业流程题（约 13 分）
  \item 实验探究题（约 13 分）
  \item 化学反应原理题（约 11 分）
\end{itemize}

选考（二选一）：
\begin{itemize}
  \item 选修三：物质结构与性质（15 分）
  \item 选修五：有机化学基础（15 分）
\end{itemize}

\subsection{答题顺序策略}

\textbf{推荐答题顺序}：
\begin{enumerate}
  \item 选择题 1--10 题（基础题，15 分钟，确保正确率 90\%+）
  \item 选考题（15 分钟，分值高且相对独立）
  \item 必考非选择题中较熟悉的一题（10 分钟）
  \item 选择题 11--14 题（中档偏难，8 分钟）
  \item 剩余必考非选择题（20 分钟，先易后难）
  \item 检查时间（7 分钟）
\end{enumerate}

\textbf{底线}：选择题目标 36+/42 分，非选择题目标 40+/58 分，总分 76+ 分。

\subsection{选择题答题技巧}

\begin{enumerate}
  \item \textbf{直接判断法}：概念题、元素推断题直接根据所学知识判断
  \item \textbf{排除法}：明显错误选项先排除
  \item \textbf{比较法}：对比选项差异，结合化学原理分析
  \item \textbf{守恒法}：电子守恒、电荷守恒、原子守恒快速判断
  \item \textbf{特殊值法}：代入特殊元素或特殊情况检验
\end{enumerate}

\textbf{选择题常见陷阱}：
\begin{itemize}
  \item 阿伏加德罗常数题：注意标准状况、气体状态、弱电解质电离、可逆反应
  \item 离子方程式正误：检查拆分是否正确、是否符合事实、是否守恒
  \item 周期律推断：找准突破口（原子序数、原子半径、化合价）
  \item 实验操作：注意"量"的问题（过量/少量）、试剂顺序
  \item 有机题：注意同分异构体计数、官能团性质判断
\end{itemize}

\subsection{工业流程题答题技巧}

\textbf{常见设问方式}：
\begin{itemize}
  \item 写出某步反应的化学/离子方程式
  \item 分析某步操作的目的（除杂、富集、调节 pH 等）
  \item 判断滤渣/滤液的成分
  \item 计算产率或纯度
  \item 分析循环利用的物质
\end{itemize}

\textbf{解题步骤}：
\begin{enumerate}
  \item 读题干：明确原料和目标产物
  \item 分步骤：分析每一步的化学反应和操作目的
  \item 看流程：关注进出物质，找出循环和副产品
  \item 答问题：结合化学原理逐空填写
\end{enumerate}

\textbf{常用规律}：
\begin{itemize}
  \item 调节 pH 的目的是使某种金属离子沉淀
  \item "趁热过滤"防止某物质降温析出
  \item "洗涤"的目的：除去表面可溶性杂质
  \item "灼烧/煅烧"发生分解反应
\end{itemize}

\subsection{实验探究题答题技巧}

\textbf{常见题型}：
\begin{itemize}
  \item 气体制备与性质验证（装置连接）
  \item 定量实验（中和滴定、重量分析）
  \item 探究性实验（提出假设 $\to$ 设计实验 $\to$ 得出结论）
  \item 物质检验与鉴别
\end{itemize}

\textbf{解题要点}：
\begin{enumerate}
  \item 明确实验目的（验证什么性质、测定什么数据）
  \item 装置连接顺序：发生 $\to$ 除杂 $\to$ 干燥 $\to$ 性质检验 $\to$ 尾气处理
  \item 注意：先除杂后干燥、防倒吸、防堵塞
  \item 现象描述要完整（"什么 + 在哪儿 + 怎么样"）
  \item 误差分析：从原理和操作两方面分析
\end{enumerate}

\textbf{常用答题模板}：
\begin{itemize}
  \item 检查装置气密性：关闭开关，微热，导管口有气泡，冷却后有水柱
  \item 作用：液封（防漏气）、平衡气压（防堵塞）
  \item 尾气处理：酸性气体用碱液，$\text{CO}$ 用燃烧法
\end{itemize}

\subsection{化学反应原理题答题技巧}

\textbf{常见考点}：
\begin{itemize}
  \item 热化学方程式书写（盖斯定律）
  \item 化学反应速率计算与影响因素
  \item 化学平衡移动分析
  \item 平衡常数计算与转化率
  \item 电离平衡、水解平衡、沉淀溶解平衡
  \item 电化学（原电池/电解池）
\end{itemize}

\textbf{解题模板}：
\begin{enumerate}
  \item 热化学：先写方程式再标状态和 $\Delta H$
  \item 速率与平衡：用\"三段式\"计算
  \item 平衡移动：勒夏特列原理 + 浓度商（$Q_c$）与 $K$ 比较
  \item 电化学：先判断电极（负/正、阴/阳）再写电极反应
  \item 溶液中的离子平衡：三大守恒（电荷、物料、质子）\footnote{电荷守恒：溶液中阳离子所带正电荷总数等于阴离子所带负电荷总数；物料守恒：某元素原子总数在溶液中等于其在各种粒子中的总数之和；质子守恒：得质子数等于失质子数。}
\end{enumerate}

\textbf{易错提醒}：
\begin{itemize}
  \item 书写 $K$ 表达式时固体和纯溶剂不写入
  \item $\Delta H$ 单位是 $\text{kJ/mol}$，注意符号和数值
  \item 原电池中阳离子向正极移动，电解池中阳离子向阴极移动
\end{itemize}

\subsection{有机推断题答题技巧}

\textbf{解题步骤}：
\begin{enumerate}
  \item 确定分子式（质谱、元素分析）
  \item 计算不饱和度（$\Omega = \frac{2n_C+2-n_H}{2}$，含 N 同理）
  \item 推断官能团（特征反应 $\to$ 可能结构）
  \item 碳链结构（通过反应条件、转化关系推断）
  \item 写出结构简式
  \item 核验（分子式、性质、反应条件是否吻合）
\end{enumerate}

\textbf{常见转化关系}：

\[
\begin{array}{c}
\text{卤代烃} \xrightarrow[\Delta]{\text{NaOH/H}_2\text{O}} \text{醇} \xrightarrow[\text{浓}\text{H}_2\text{SO}_4]{\Delta} \text{烯烃}\\
\text{醇} \xrightarrow[\Delta]{\text{Cu/O}_2} \text{醛} \xrightarrow[\Delta]{\text{Cu(OH)}_2} \text{羧酸}\\
\text{醛} \xrightarrow{\text{H}_2/\text{Ni}} \text{醇}\\
\text{羧酸} + \text{醇} \xrightarrow[\Delta]{\text{浓}\text{H}_2\text{SO}_4} \text{酯} + \text{H}_2\text{O}\\
\text{酯} \xrightarrow[\Delta]{\text{H}_2\text{O}/\text{H}^+} \text{羧酸} + \text{醇}\\
\text{酯} \xrightarrow[\Delta]{\text{H}_2\text{O}/\text{OH}^-} \text{羧酸盐} + \text{醇}
\end{array}
\]

\textbf{易错提醒}：
\begin{itemize}
  \item 卤代烃消去反应条件：$\text{NaOH}/\text{乙醇}$，加热
  \item 醇消去：浓 $\text{H}_2\text{SO}_4$、$170^\circ\text{C}$
  \item 酚酯水解消耗 2 倍 $\text{NaOH}$
  \item $\text{Fe/HCl}$ 还原硝基为氨基
\end{itemize}

\subsection{物质结构与性质题答题技巧}

\textbf{常见考点}：
\begin{itemize}
  \item 电子排布式/轨道表示式
  \item 第一电离能、电负性比较
  \item 杂化轨道类型与空间构型
  \item 分子间作用力（氢键）与熔沸点
  \item 晶体类型与晶胞计算
\end{itemize}

\textbf{晶胞计算常用公式}：
\[
\rho = \frac{Z \cdot M}{N_A \cdot a^3}
\]
（$Z$：晶胞中微粒数，$M$：摩尔质量，$a$：晶胞棱长）

\textbf{空间利用率}：
\[
\text{体心立方}：\frac{\sqrt{3}\pi}{8} \approx 68\%,\qquad
\text{面心立方}：\frac{\sqrt{2}\pi}{6} \approx 74\%
\]

\textbf{易错提醒}：
\begin{itemize}
  \item 书写电子排布式注意能级交错：$4s < 3d$ 但失电子先失 $4s$
  \item 比较电离能时注意全满/半满规则（如 $\text{N} > \text{O}$）
  \item 氢键只影响物理性质（熔沸点、溶解度），不影响化学性质
\end{itemize}

\subsection{考场应急策略}

\begin{enumerate}
  \item \textbf{遇到难题}：超过 2 分钟无思路先跳过，标记后回头再做
  \item \textbf{计算卡壳}：检查单位换算、数据代入是否正确
  \item \textbf{有机推断卡壳}：从不饱和度 $\to$ 可能结构 $\to$ 反推验证
  \item \textbf{实验题不明}：回到实验目的，分析每一步操作的目的
  \item \textbf{剩余 5 分钟}：检查选择题涂卡、非选择题是否漏空
  \item \textbf{剩余 1 分钟}：确保选考涂卡标记正确
\end{enumerate}
